% !PRINTED 448
\weberpage{448}
%<<CONT
\weq{22}{\begin{aligned}
(k)^p &\equiv (kl) & &(\mathrm{mod}\ \mathfrak P),\\
(kl)^{p'} &\equiv (kll') & &(\mathrm{mod}\ \mathfrak P'),\\
(kll')^{p''} &\equiv (kll'l'') & &(\mathrm{mod}\ \mathfrak P'')
\end{aligned}}
\[
\cdot\ \cdot\ \cdot\ \cdot\ \cdot\ \cdot\ \cdot\ \cdot\ \cdot\ \cdot\ \cdot
\]
Nehmen wir jetzt an, es zerfalle die Klassengleichung, es sei
also für ein variables $t$
\[
H_m(t) = H_1(t) H_2(t),
\]
so können wir, da $k, k'$ beliebige Klassen waren, $(k)$ unter den
Wurzeln von $H_1(t)$, $(k')$ unter denen von $H_2(t)$ wählen. In der
Kette (22) gehört also das Anfangsglied $(k)$ zu den Wurzeln $(k_1)$
von $H_1(t)$ und das Endglied $(k')$ zu den Wurzeln $(k_2)$ von $H_2(t)$.
Da mindestens an einer Stelle der Kette (22) der Übergang von
den Wurzeln des einen Faktors zu denen des anderen stattfinden
muß, so gibt es ein Paar von Wurzeln $(k_1), (k_2)$, das für irgend
eine Primzahl $p$ und ein darin aufgehendes Primideal $\mathfrak P$ der
Kongruenz
\[
(k_1)^p \equiv (k_2) \quad (\mathrm{mod}\ \mathfrak P)
\]
genügt. Da nun $H_1(k_1) = 0$ ist, so folgt durch den oft angewandten Schluß
\[
H_1(k_1)^p \equiv H_1(k_1^p) \equiv H_1(k_2) \equiv 0 \quad (\mathrm{mod}\ \mathfrak P).
\]
Da nun
\[
H_1(k_2) = \prod^{k_1} [(k_2) - (k_1)]
\]
ist, so muß eine der Differenzen $(k_2) - (k_1)$ durch $\mathfrak P$ teilbar sein,
was nicht möglich ist, da $p$ nicht in der Diskriminante von $H$
aufgeht.

\gesperrt{Damit ist die Irreducibilität der Klassengleichung bewiesen; der Grad des Körpers $\mathfrak K$ ist gleich
dem Doppelten der Klassenzahl $h$ festgestellt, und
die im vorigen Paragraphen gefundene Gruppe $\mathfrak S$
ist als die wahre Gruppe der Klassengleichung erkannt.}

8. Hiernach können wir alle Primzahlen $p$, \gesperrt{abgesehen
von einer endlichen Anzahl von Ausnahmen} (die in $2D$
oder der Diskriminante von $H$ aufgehen) in ihre Primfaktoren
im Körper $\mathfrak K$ zerlegen.

Ist $(D, p) = +1$, so zerfällt $p$ und Körper $\Omega$ % [sic] printed "und Körper"; grammatically and mathematically expected "im Körper"
in zwei konjugierte Primideale ersten Grades. Ist $p$ vom Index $\varepsilon$, so ist
jedes in $p$ aufgehende Primideal $\mathfrak p$ vom Grade $\varepsilon$. Ist also
\[
p = \mathfrak P_1 \mathfrak P_2 \ldots \mathfrak P_r,
\]
%CONT>>
